% apendice.tex: Apêndice A: Matriz Completa de Parâmetros Experimentais

\appendix
\chapter{Matriz de Parâmetros Experimentais}
\label{app:reproducibility}
\lhead{APÊNDICE A}

\renewcommand{\thetable}{A.\arabic{table}}
\setcounter{table}{0}

Este apêndice consolida os parâmetros utilizados nos experimentos reportados, organizados por camada experimental. A formulação conceitual de cada política permanece nos Capítulos~2 e~3; o registro aqui serve para auditoria e reprodução. Os valores foram obtidos a partir das configurações, scripts e artefatos versionados do projeto. A organização segue a sequência experimental: parâmetros globais (Tabela~\ref{tab:global_params}), políticas clássicas e meta-heurísticas (Tabela~\ref{tab:classic_meta_params}), agentes de aprendizado por reforço e políticas híbridas (Tabela~\ref{tab:rl_params}), previsão e geração de rótulos (Tabela~\ref{tab:forecast_label_params}) e ambiente de validação e reprodução (Tabela~\ref{tab:validation_env}).

% ─────────────────────────────────────────────────────────────────────────────
% TABELA A.1: Parâmetros globais do experimento
% ─────────────────────────────────────────────────────────────────────────────

\begingroup\setlength{\tabcolsep}{4pt}\scriptsize
\begin{longtable}{@{} p{2.6cm} p{5.2cm} p{2.8cm} p{4.0cm} @{}}
\caption[Parâmetros globais do experimento.]{Parâmetros globais do experimento.}
\label{tab:global_params} \\
\toprule
\textbf{Componente} & \textbf{Parâmetro} & \textbf{Valor usado} & \textbf{Observação} \\
\midrule
\endfirsthead
\multicolumn{4}{l}{\scriptsize\textit{(continuação da Tabela~A.1)}} \\[2pt]
\toprule
\textbf{Componente} & \textbf{Parâmetro} & \textbf{Valor usado} & \textbf{Observação} \\
\midrule
\endhead
\midrule
\multicolumn{4}{r}{\scriptsize\textit{(continua)}} \\
\endfoot
\bottomrule
\endlastfoot

\multicolumn{4}{@{}l}{\textit{Dados}} \\[2pt]
Dados & estado & BA & UF selecionada para a análise \\
Dados & período & jan/2023 a abr/2025 & 17 ciclos por ano \\
Dados & registros transacionais & ${\approx}48$ milhões & base completa da rede varejista \\
Dados & produtos & ${\approx}12\,000$ & base completa \\
Dados & periodicidade & ciclos sazonais & 17 ciclos por ano \\
Dados & tratamento de zeros & preenche com zero & ciclos ausentes completados \\
Dados & filtro min.\ de ciclos com venda & 17 & séries com menos ciclos positivos são removidas \\
Dados & janela de atividade do produto & 3 ciclos & produto considerado ativo se vendeu nos últimos 3 ciclos \\
Dados & status incluídos & Ativo & produtos com status diferente são excluídos \\
Dados & tipos de venda excluídos & Brinde & transações de brinde não entram na simulação \\
Dados & deduplicação & por armazém, produto, loja e ciclo & chave de deduplicação dos registros \\[4pt]

\multicolumn{4}{@{}l}{\textit{Ambiente de simulação}} \\[2pt]
InventoryEnv & lead time & 2 ciclos & uniforme; mesmo valor para todas as políticas \\
InventoryEnv & estoque inicial $I_0$ & escalonado por série & $\mathrm{rop\_ref} = \mu L + z\sigma\sqrt{L}$; ajustado por série \\
InventoryEnv & tipo de falta & perda de venda & demanda não atendida é perdida; sem acúmulo \\
InventoryEnv & regra de arredondamento & inteiro superior & aplicada a todos os pedidos calculados \\
InventoryEnv & frequência de reavaliação (AIPE) & 6 ciclos & intervalo $\Delta t$ definido conceitualmente \\[4pt]

\multicolumn{4}{@{}l}{\textit{Custos}} \\[2pt]
Custos & custo de manutenção $h$ & 1,0 & por unidade por ciclo \\
Custos & custo de ruptura $c_s$ & 5,0 & por unidade de demanda não atendida \\
Custos & custo fixo de pedido $K$ & 50,0 & por ordem emitida \\
Custos & custo variável de pedido $c_v$ & 0,5 & por unidade pedida \\
Custos & razão $c_s/h$ & 5,0 & nível de serviço crítico ${\approx}83\%$ \\
Custos & unidade monetária & adimensional & valores relativos, não em R\$ \\
Custos & custos globais ou por série & globais & mesmos parâmetros para todas as séries \\[4pt]

\multicolumn{4}{@{}l}{\textit{Avaliação}} \\[2pt]
Avaliação & número de replicações & 5 & replicações independentes por série \\
Avaliação & semente global & 42 & fixa para todas as execuções \\
Avaliação & modo de avaliação & \textit{walk-forward} & janela deslizante; 17 ciclos de treino \\
Avaliação & métricas reportadas & CTI, NS, TR, BE, FP & definição formal no Capítulo~3 \\
Avaliação & limiar mínimo de NS (viabilidade) & 0,70 & política abaixo do limiar é considerada inviável \\
\end{longtable}
\endgroup

% ─────────────────────────────────────────────────────────────────────────────
% TABELA A.2: Políticas clássicas e meta-heurísticas
% ─────────────────────────────────────────────────────────────────────────────

Nas políticas clássicas, os parâmetros de reposição são calculados diretamente a partir das estatísticas históricas da série e dos parâmetros globais de custo. Nas meta-heurísticas, por sua vez, cada candidato é uma política $(s,Q)$, em que $s$ representa o ponto de reposição e $Q$ representa a quantidade de pedido. A avaliação de cada candidato ocorre pela execução da política correspondente no ambiente de simulação, retornando os indicadores de custo total de inventário, nível de serviço, ruptura, efeito chicote e frequência de pedidos.

Para o Algoritmo Genético, cada indivíduo da população é codificado como um vetor contínuo $(s,Q)$. A população inicial é amostrada dentro dos limites operacionais admissíveis para a série avaliada. A seleção é realizada por torneio, favorecendo indivíduos com maior aptidão. O operador de cruzamento combina os parâmetros de dois indivíduos pais por meio de cruzamento do tipo \textit{Blend}, enquanto a mutação aplica perturbação gaussiana aos genes, seguida de ajuste aos limites admissíveis. Dessa forma, o GA busca combinações de $(s,Q)$ que maximizem a aptidão definida para o experimento.

As demais meta-heurísticas usam a mesma representação de solução e a mesma função de avaliação. No SA, uma solução inicial $(s,Q)$ é iterativamente perturbada por movimentos de vizinhança, aceitando pioras de custo com probabilidade controlada pela temperatura e pelo cronograma de resfriamento. No PSO, cada partícula representa um par $(s,Q)$ com velocidade própria, sendo atualizada por termos de inércia, atração ao melhor ponto individual e atração ao melhor ponto global do enxame. No DE, cada vetor $(s,Q)$ é atualizado por mutação diferencial e cruzamento binomial, preservando apenas candidatos que melhoram a função objetivo.

\begingroup\setlength{\tabcolsep}{4pt}\scriptsize
\begin{longtable}{@{} p{2.0cm} p{5.5cm} p{3.0cm} p{4.1cm} @{}}
\caption[Parâmetros das políticas clássicas e meta-heurísticas.]{Parâmetros das políticas clássicas e meta-heurísticas.}
\label{tab:classic_meta_params} \\
\toprule
\textbf{Política} & \textbf{Parâmetro} & \textbf{Valor usado} & \textbf{Observação} \\
\midrule
\endfirsthead
\multicolumn{4}{l}{\scriptsize\textit{(continuação da Tabela~A.2)}} \\[2pt]
\toprule
\textbf{Política} & \textbf{Parâmetro} & \textbf{Valor usado} & \textbf{Observação} \\
\midrule
\endhead
\midrule
\multicolumn{4}{r}{\scriptsize\textit{(continua)}} \\
\endfoot
\bottomrule
\endlastfoot

\multicolumn{4}{@{}l}{\textit{EOQ}} \\[2pt]
EOQ & lote ótimo $Q^*$ & $\sqrt{2KD/h}$ & $K{=}50$, $h{=}1$ (globais) \\
EOQ & estimativa de demanda $D$ & previsão auxiliar & série histórica filtrada \\
EOQ & ponto de reposição & $\mu L + z\sigma\sqrt{L}$ & $L{=}2$; $z{\approx}0{,}97$ \\[4pt]

\multicolumn{4}{@{}l}{\textit{$(s,S)$}} \\[2pt]
$(s,S)$ & ponto de disparo $s$ & ROP $= \mu L + z\sigma\sqrt{L}$ & $L{=}2$ \\
$(s,S)$ & nível alvo $S$ & $s + Q^*$ & derivado do EOQ \\
$(s,S)$ & fator de segurança $z$ & $\Phi^{-1}(c_s/(c_s{+}h)) \approx 0{,}97$ & $c_s{=}5$, $h{=}1$ \\[4pt]

\multicolumn{4}{@{}l}{\textit{Jornaleiro (Newsvendor)}} \\[2pt]
Jornaleiro & razão crítica CR & $c_s/(c_s{+}h) = 5/6 \approx 0{,}833$ & \\
Jornaleiro & nível ótimo $y^*$ & quantil(0,833) da previsão & via distribuição empírica \\[4pt]

\multicolumn{4}{@{}l}{\textit{GA (Algoritmo Genético)}} \\[2pt]
GA & variáveis otimizadas & ponto de reposição e lote & $(s, Q)$ contínuos \\
GA & tamanho da população & 100 & \\
GA & número de gerações & 50 & \\
GA & taxa de crossover & 0,9 & operador Blend, $\alpha{=}0{,}5$ \\
GA & taxa de mutação & 0,05 & gaussiana, $\sigma{=}\mathrm{range}/10$ \\
GA & método de seleção & torneio de tamanho 3 & \\
GA & função fitness & $1{,}0 \times \mathrm{NS} - 0{,}0001 \times \mathrm{CTI}$ & maximiza NS com penalidade de custo \\
GA & semente & 42 & \\[4pt]

\multicolumn{4}{@{}l}{\textit{SA (Simulated Annealing)}} \\[2pt]
SA & variáveis otimizadas & ponto de reposição e lote & $(s, Q)$ contínuos \\
SA & temperatura inicial & 1000,0 & \\
SA & fator de resfriamento & 0,95 & geométrico por iteração \\
SA & iterações máximas & 500 & \\
SA & minimizador local & Nelder-Mead & \texttt{scipy.optimize.}\newline\texttt{dual\_annealing} \\
SA & semente & 42 & \\[4pt]

\multicolumn{4}{@{}l}{\textit{PSO (Particle Swarm Optimization)}} \\[2pt]
PSO & variáveis otimizadas & ponto de reposição e lote & $(s, Q)$ contínuos \\
PSO & número de partículas & 40 & \\
PSO & iterações máximas & 80 & \\
PSO & peso de inércia $w$ & 0,7 & \\
PSO & coeficiente cognitivo $c_1$ & 1,5 & \\
PSO & coeficiente social $c_2$ & 1,5 & \\
PSO & semente & 42 & \\[4pt]

\multicolumn{4}{@{}l}{\textit{DE (Differential Evolution)}} \\[2pt]
DE & variáveis otimizadas & ponto de reposição e lote & $(s, Q)$ contínuos \\
DE & estratégia & \texttt{rand1bin} & mutação + cruzamento binomial \\
DE & fator diferencial $F$ & 0,8 & \\
DE & taxa de cruzamento CR & 0,9 & \\
DE & iterações máximas & 100 & \\
DE & tamanho da população & 15 & \\
DE & polimento local & ativo & \texttt{scipy.optimize.}\newline\texttt{differential\_evolution} \\
DE & tolerância & $10^{-6}$ & critério de parada \\
DE & semente & 42 & \\
\end{longtable}
\endgroup

% ─────────────────────────────────────────────────────────────────────────────
% TABELA A.3: Agentes RL e políticas híbridas
% ─────────────────────────────────────────────────────────────────────────────

Nos agentes de aprendizado por reforço, a decisão de reposição é formulada como um problema sequencial no ambiente \textit{InventoryEnv}. O estado observado pelo agente resume a posição de estoque, demanda prevista, pedidos em trânsito, estatísticas recentes de demanda e progresso no horizonte experimental. A ação representa a quantidade de pedido, discreta no caso do DQN e do SARSA e contínua no caso do PPO. A recompensa é definida como o negativo do custo operacional do ciclo, de modo que maximizar retorno equivale a reduzir custo total, penalizando simultaneamente estoque, ruptura e emissão de pedidos.

Nas políticas híbridas GA-DQN e GA-PPO, o GA fornece uma política inicial de boa qualidade para reduzir o custo de exploração dos agentes de RL em horizontes curtos. No GA-DQN, trajetórias geradas pela política calibrada pelo GA pré-populam o \textit{replay buffer}, e o agente refina a decisão por episódios adicionais. No GA-PPO, os \textit{rollouts} iniciais produzidos pelo GA servem como aquecimento antes do refinamento da política contínua. Em ambos os casos, a comparação permanece controlada porque o ambiente, os custos, o horizonte e as séries de demanda são os mesmos das demais políticas.

\begingroup\setlength{\tabcolsep}{4pt}\scriptsize
\begin{longtable}{@{} p{2.0cm} p{5.5cm} p{3.2cm} p{3.9cm} @{}}
\caption[Parâmetros dos agentes de aprendizado por reforço e políticas híbridas.]{Parâmetros dos agentes de aprendizado por reforço e políticas híbridas.}
\label{tab:rl_params} \\
\toprule
\textbf{Política} & \textbf{Parâmetro} & \textbf{Valor usado} & \textbf{Observação} \\
\midrule
\endfirsthead
\multicolumn{4}{l}{\scriptsize\textit{(continuação da Tabela~A.3)}} \\[2pt]
\toprule
\textbf{Política} & \textbf{Parâmetro} & \textbf{Valor usado} & \textbf{Observação} \\
\midrule
\endhead
\midrule
\multicolumn{4}{r}{\scriptsize\textit{(continua)}} \\
\endfoot
\bottomrule
\endlastfoot

\multicolumn{4}{@{}l}{\textit{Ambiente compartilhado (todos os agentes)}} \\[2pt]
InventoryEnv & vetor de estado (dim.\ 1 a 3) & $I_t/I_0$,\; $\hat{D}/\mu$,\; $\mathrm{OP}/I_0$ & nível de estoque, previsão, pedido aberto \\
InventoryEnv & vetor de estado (dim.\ 4 a 6) & $\mu_6/\mu$,\; $\sigma_6/\sigma$,\; $t/T$ & tendências de curto prazo e posição no horizonte \\
InventoryEnv & função de recompensa & $-(h I_t + c_s S_t + K\mathbf{1}_{Q>0} + c_v Q_t)$ & $S_t$ = demanda não atendida \\[4pt]

\multicolumn{4}{@{}l}{\textit{DQN}} \\[2pt]
DQN & espaço de ações & 20 ações discretas & $Q \in [0,\, \max(200,\, 12\mu)]$ \\
DQN & número de episódios & 500 & \\
DQN & fator de desconto $\gamma$ & 0,95 & \\
DQN & $\varepsilon$ inicial & 1,0 & exploração total no início \\
DQN & $\varepsilon$ mínimo & 0,01 & \\
DQN & decaimento de $\varepsilon$ & 0,995 por episódio & geométrico \\
DQN & tamanho do batch & 64 & \\
DQN & tamanho do \textit{replay buffer} & 10\,000 & \\
DQN & frequência de atualização da rede-alvo & 10 episódios & \\
DQN & arquitetura da Q-network & $[6] \to [64, 64] \to [20]$ & duas camadas ocultas \\
DQN & taxa de aprendizado & 0,001 & padrão do código; não exposto no YAML \\
DQN & recorte de gradiente & $[-1{,}0,\; 1{,}0]$ & \\
DQN & semente & 42 & \\[4pt]

\multicolumn{4}{@{}l}{\textit{PPO}} \\[2pt]
PPO & espaço de ações & 20 ações discretas & $Q \in [0,\, \max(200,\, 12\mu)]$ \\
PPO & número de episódios & 500 & \\
PPO & fator de desconto $\gamma$ & 0,99 & \\
PPO & epsilon de recorte & 0,2 & razão de probabilidade recortada em $[1{-}\varepsilon,\, 1{+}\varepsilon]$ \\
PPO & épocas de otimização por rollout & 4 & \\
PPO & taxa de aprendizado & 0,0003 & \\
PPO & arquitetura do ator & $[6] \to [64, 32] \to [20]$ & saída em logits \\
PPO & arquitetura do crítico & $[6] \to [64, 32] \to [1]$ & saída em valor escalar \\
PPO & estimativa de vantagem & $A_t = G_t - V(s_t)$, normalizada & $(G - \bar{G})/(\sigma_G + 10^{-8})$ \\
PPO & semente & 42 & \\[4pt]

\multicolumn{4}{@{}l}{\textit{SARSA}} \\[2pt]
SARSA & discretização do estado & $\lfloor \mathrm{clip}(I_t/I_0,\, 0,\, 1) \times 5 \rfloor$ & bin único baseado no nível de estoque \\
SARSA & número de estados & 5 & \\
SARSA & espaço de ações & 10 ações discretas & \\
SARSA & política de exploração & $\varepsilon$-greedy & \\
SARSA & número de episódios & 500 & \\
SARSA & fator de desconto $\gamma$ & 0,99 & \\
SARSA & taxa de aprendizado $\alpha$ & 0,1 & \\
SARSA & $\varepsilon$ & 0,1 & constante; sem decaimento identificado \\
SARSA & semente & 42 & \\[4pt]

\multicolumn{4}{@{}l}{\textit{GA-DQN}} \\[2pt]
GA-DQN & parâmetros do GA & herdados (ver Tabela~A.2) & 50 gerações \\
GA-DQN & parâmetros do DQN & herdados (ver acima) & \\
GA-DQN & tamanho do \textit{replay buffer} pré-populado & 1\,000 & trajetórias geradas pelo GA \\
GA-DQN & episódios de refinamento RL & 200 & \\
GA-DQN & regra de integração GA~$\to$~DQN & se $(I_t + \mathrm{OP}) \leq s + \mathrm{SS}$:\newline $Q = \max(Q_{\mathrm{DQN}},\; 0{,}5 \cdot Q_{\mathrm{GA}})$;\newline caso contrário: $Q = 0$ & \\
GA-DQN & semente & 42 & \\[4pt]

\multicolumn{4}{@{}l}{\textit{GA-PPO}} \\[2pt]
GA-PPO & parâmetros do GA & herdados (ver Tabela~A.2) & 50 gerações \\
GA-PPO & parâmetros do PPO & herdados (ver acima) & \\
GA-PPO & episódios de aquecimento com trajetórias GA & 20 & log-prob uniforme \\
GA-PPO & episódios de refinamento PPO & 200 & \\
GA-PPO & forma de integração GA~$\to$~PPO & inicialização por rollouts GA; refinamento PPO & warm-start antes do treinamento RL \\
GA-PPO & semente & 42 & \\
\end{longtable}
\endgroup

% ─────────────────────────────────────────────────────────────────────────────
% TABELA A.4: Previsão, características operacionais e geração de rótulos
% ─────────────────────────────────────────────────────────────────────────────

\begingroup\setlength{\tabcolsep}{4pt}\scriptsize
\begin{longtable}{@{} p{2.6cm} p{5.2cm} p{2.8cm} p{4.0cm} @{}}
\caption[Parâmetros de previsão, características operacionais e geração de rótulos.]{Parâmetros de previsão, características operacionais e geração de rótulos.}
\label{tab:forecast_label_params} \\
\toprule
\textbf{Componente} & \textbf{Parâmetro} & \textbf{Valor usado} & \textbf{Observação} \\
\midrule
\endfirsthead
\multicolumn{4}{l}{\scriptsize\textit{(continuação da Tabela~A.4)}} \\[2pt]
\toprule
\textbf{Componente} & \textbf{Parâmetro} & \textbf{Valor usado} & \textbf{Observação} \\
\midrule
\endhead
\midrule
\multicolumn{4}{r}{\scriptsize\textit{(continua)}} \\
\endfoot
\bottomrule
\endlastfoot

\multicolumn{4}{@{}l}{\textit{Previsão auxiliar}} \\[2pt]
Previsão & modelos utilizados & LSTM, ANN, XGBoost, Croston, ARIMA & ensemble de previsão \\
Previsão & janela de entrada (lookback) & 6 ciclos & \\
Previsão & ciclos mínimos de treino & 17 & \\
Previsão & modo \textit{walk-forward} & desativado & previsão em janela fixa \\
Previsão & horizonte de previsão & 1 ciclo à frente & \\[3pt]

LSTM & hidden size & 64 & \\
LSTM & épocas & 50 & \\
LSTM & batch size & 16 & \\
LSTM & taxa de aprendizado & 0,001 & \\[3pt]

ANN / MLP & camadas ocultas & $[64, 32]$ & \\
ANN / MLP & iterações máximas & 200 & \\
ANN / MLP & early stopping & ativo & \\
ANN / MLP & fração de validação & 0,1 & \\[3pt]

XGBoost (previsão) & n\_estimators & 200 & \\
XGBoost (previsão) & max\_depth & 6 & \\
XGBoost (previsão) & taxa de aprendizado & 0,05 & \\
XGBoost (previsão) & \textit{subsample} & 0,8 & \\[3pt]

Croston & alpha (suavização) & 0,1 & suavização exponencial \\[3pt]

ARIMA & ordem $(p,d,q)$ & $(1,1,1)$ & \\[4pt]

\multicolumn{4}{@{}l}{\textit{Características operacionais e classificação POD}} \\[2pt]
POD / ADI-CV$^2$ & limiar ADI & 1,32 & classif.\ Syntetos-Boylan \\
POD / ADI-CV$^2$ & limiar CV$^2$ & 0,49 & classif.\ Syntetos-Boylan \\
POD / ADI-CV$^2$ & limiar de zeros consecutivos & 5 & \\
POD / ADI-CV$^2$ & limiar de burstiness & 0,30 & \\
POD / ADI-CV$^2$ & limiar ACF (sazonalidade) & 0,25 & \\
POD / ADI-CV$^2$ & coeficiente mínimo de tendência & 0,05 & \\
POD / ADI-CV$^2$ & limiar de entropia & 2,0 & \\
POD / ADI-CV$^2$ & média mínima de demanda & 2,0 & séries abaixo são excluídas \\
POD / ADI-CV$^2$ & limiar CV para segmentação & 1,5 & segmentação adicional por variabilidade \\[4pt]

\multicolumn{4}{@{}l}{\textit{Geração de rótulos e PSE}} \\[2pt]
Rótulos & limiar mínimo de NS para rótulo & 0,70 & política abaixo é excluída do \textit{ranking} \\
Rótulos & modelo PSE & XGBoost & classificador treinado sobre perfis de séries \\
Rótulos & validação cruzada & 5 folds & \\
Rótulos & semente PSE & 42 & \\
Rótulos & n\_estimators (PSE) & 300 & \\
Rótulos & max\_depth (PSE) & 6 & \\
Rótulos & taxa de aprendizado (PSE) & 0,05 & \\
Rótulos & \textit{subsample} (PSE) & 0,8 & \\
Rótulos & colsample\_bytree (PSE) & 0,8 & \\
\end{longtable}
\endgroup

% ─────────────────────────────────────────────────────────────────────────────
% TABELA A.5: Validação estatística, ambiente computacional e artefatos
% ─────────────────────────────────────────────────────────────────────────────

\begingroup\setlength{\tabcolsep}{4pt}\scriptsize
\begin{longtable}{@{} p{2.6cm} p{5.2cm} p{2.8cm} p{4.0cm} @{}}
\caption[Validação estatística, ambiente computacional e artefatos.]{Validação estatística, ambiente computacional e artefatos.}
\label{tab:validation_env} \\
\toprule
\textbf{Componente} & \textbf{Parâmetro ou artefato} & \textbf{Valor usado} & \textbf{Observação} \\
\midrule
\endfirsthead
\multicolumn{4}{l}{\scriptsize\textit{(continuação da Tabela~A.5)}} \\[2pt]
\toprule
\textbf{Componente} & \textbf{Parâmetro ou artefato} & \textbf{Valor usado} & \textbf{Observação} \\
\midrule
\endhead
\midrule
\multicolumn{4}{r}{\scriptsize\textit{(continua)}} \\
\endfoot
\bottomrule
\endlastfoot

\multicolumn{4}{@{}l}{\textit{Validação estatística}} \\[2pt]
Estatística & teste par a par & Wilcoxon signed-rank & comparação de pares de políticas \\
Estatística & teste global & Friedman & hipótese nula de igualdade entre múltiplas políticas \\
Estatística & \textit{post-hoc} & Nemenyi & comparações múltiplas após rejeição de Friedman \\[4pt]

\multicolumn{4}{@{}l}{\textit{Ambiente computacional}} \\[2pt]
Python & versão mínima & 3.9 & \\
Kedro & versão mínima & 0.19.0 & orquestração do \textit{pipeline} \\
kedro-datasets & versão mínima & 2.0.0 & \\
pandas & versão mínima & 2.0 & \\
NumPy & versão mínima & 1.24 & \\
scikit-learn & versão mínima & 1.3 & \\
XGBoost & versão mínima & 2.0 & \\
DEAP & versão mínima & 1.4 & GA e DE \\
SciPy & versão mínima & 1.11 & SA e DE \\
PyArrow & versão mínima & 14.0 & \\
Matplotlib & versão mínima & 3.7 & \\
scikit-posthocs & versão mínima & 0.8 & testes de Nemenyi \\
GeoPandas & versão mínima & 0.14 & \\
geobr & versão mínima & 0.4 & malhas geográficas \\
PyYAML & versão mínima & 6.0 & \\[4pt]

\multicolumn{4}{@{}l}{\textit{Artefatos de reprodução}} \\[2pt]
dados processados & \nolinkurl{simulation/data/03_primary/} & & snapshot dos dados usados nas simulações \\
características operacionais & \nolinkurl{simulation/data/04_feature/} & & descritores das séries de demanda \\
configurações YAML & \nolinkurl{simulation/conf/base/parameters/} & & hiperparâmetros de todos os módulos \\
resultados de simulação & \nolinkurl{simulation/data/07_model_output/} & & métricas CTI, NS, TR, BE e FP por série \\
testes estatísticos & \nolinkurl{results/statistical_tests/} & & Wilcoxon, Friedman, Nemenyi \\
rótulos PSE & \nolinkurl{results/policy_labels/} & & política dominante por série \\
\end{longtable}
\endgroup

% ─────────────────────────────────────────────────────────────────────────────
% TABELA A.6: Resultados agregados da Fase 1 (Paraíba)
% ─────────────────────────────────────────────────────────────────────────────

\section{Resultados Agregados da Fase 1}
\label{app:fase1_resultados}

Esta seção registra a tabela completa de indicadores agregados da Fase~1 (15 lojas da Paraíba, produto~48130, regime \textit{Lumpy}), citada na Seção~\ref{sec:custo_politica_unica_vs_perfil} apenas em forma resumida para evitar duplicação com a comparação agregada da Fase~2 (Tabela~\ref{tab:kpis_fase2}).

\noindent\textbf{Como ler a tabela:} políticas com melhor equilíbrio combinam baixo CTI, alto NS, baixa TR e baixo BE. FP indica a frequência de emissão de pedidos: valores próximos de zero ($\approx 0$) ou de um ($\approx 1$) sinalizam degeneração operacional (o agente parou de emitir pedidos ou pediu em quase todos os ciclos), e essas políticas não constituem soluções práticas viáveis.

\begin{table}[ht]
\centering
\small
\caption{Médias dos indicadores de desempenho por política (15 lojas, regime \textit{Lumpy}). Melhores valores por coluna. $^\dagger$SARSA e PPO apresentam degeneração (FP\,=\,0 e FP\,=\,0.98, respectivamente); resultados degenerados não constituem soluções operacionalmente viáveis.}
\label{tab:global_prelim}
\begin{tabular}{@{}lrrrrr@{}}
\toprule
Política &
\makecell[r]{CTI\\(R\$)} &
\makecell[r]{NS} &
\makecell[r]{TR} &
\makecell[r]{BE} &
\makecell[r]{FP} \\
\midrule
SARSA$^\dagger$              & 16.139 & 0.785 & 0.200 & 0.00 & 0.00 \\
Jornaleiro                   & 16.402 & 0.794 & 0.111 & 0.003 & 0.27 \\
DQN                          & 17.974 & 0.806 & 0.078 & 4.1   & 0.37 \\
GA-DQN (Híbrido)             & 18.502 & 0.821 & 0.050 & 5.0   & 0.31 \\
PSO                          & 19.100 & 0.819 & 0.056 & 6.2   & 0.31 \\
GA-PPO (Híbrido)             & 19.801 & 0.833 & 0.050 & 6.6   & 0.35 \\
GA                           & 20.015 & 0.825 & 0.061 & 19.7  & 0.30 \\
DE                           & 20.252 & 0.813 & 0.056 & 31.0  & 0.27 \\
SA                           & 20.967 & 0.820 & 0.056 & 36.4  & 0.27 \\
EOQ                          & 31.034 & 0.840 & 0.050 & 23.0  & 0.20 \\
PPO$^\dagger$                & 34.637 & 0.848 & 0.039 & 15.9 & 0.98 \\
$(s,S)$                      & 35.589 & 0.843 & 0.050 & 119.4 & 0.15 \\
\bottomrule
\end{tabular}
\end{table}

Os dois extremos revelam os riscos da otimização por única métrica: o SARSA$^\dagger$ converge ao comportamento degenerado de não emitir pedidos (FP\,=\,0), obtendo o menor CTI ao preço de 20\% de rupturas; o PPO$^\dagger$ emite pedidos em 98\% dos ciclos, atingindo o maior NS mas com custo inviável.

\noindent A redução de ``até 48\%'' citada no resumo deste trabalho origina-se exatamente desta comparação, calculada diretamente da Tabela~\ref{tab:global_prelim}: CTI médio GA-DQN (R\$\,18.502) versus $(s,S)$ (R\$\,35.589), $(35.589-18.502)/35.589 \approx 48\%$, sobre as 15 lojas da Fase~1, produto~48130 (Paraíba). O valor representa o ganho máximo observado neste experimento, não uma redução média universal: na Fase~2 (145 séries, BA) o GA-DQN não domina o EOQ em CTI, e o Jornaleiro apresenta menor custo bruto, mas não atinge o nível de serviço mínimo de $0{,}70$ (Seção~\ref{sec:custo_politica_unica_vs_perfil}).
